\phantomsection\addcontentsline{toc}{section}{Autonomous Vehicles}
\ECchapter{23}{Autonomous Vehicles}{Engineering safe autonomous mobility}{ECOrange}

\ECObjectives{
\begin{itemize}
\item Describe perception, localisation, prediction, planning and control.
\item Explain why sensor diversity and redundancy are necessary.
\item Calculate stopping distance and identify operational limits.
\end{itemize}}

\begin{multicols}{2}
\begin{engineeringnews}
Autonomous driving systems are being tested in passenger cars, shuttles, mines, ports and warehouses.
Progress is fastest in controlled environments, whereas public roads remain difficult because weather,
road geometry, human behaviour and rare events are highly variable.
\end{engineeringnews}

\begin{ecarticle}
An autonomous vehicle must perceive its environment, estimate its own position, predict how other road
users may move, plan a safe trajectory and control steering, braking and propulsion. These functions run
continuously and must satisfy strict timing, reliability and safety requirements.

Perception combines cameras, radar, lidar and ultrasonic sensors. Cameras provide colour and detailed
shapes but can be affected by darkness, glare or fog. Radar measures distance and relative speed and
remains useful in poor visibility. Lidar produces accurate three-dimensional point clouds but may be
costly and sensitive to rain or contamination. Sensor fusion combines these strengths and highlights
inconsistent measurements.

Localisation uses satellite positioning, inertial sensors, wheel odometry and maps. No single source is
reliable everywhere: GNSS signals can be blocked in tunnels, while wheel slip creates odometry errors.
Redundant estimates allow the system to detect a faulty measurement and maintain a safe state.

The prediction module estimates the future motion of pedestrians, cyclists and vehicles. The planner
then selects a manoeuvre and generates a trajectory that respects road geometry, traffic rules, vehicle
dynamics and passenger comfort. The controller follows this trajectory by adjusting steering torque,
brake pressure and propulsion.

Braking distance increases approximately with the square of speed because kinetic energy is
proportional to $v^2$. Total stopping distance also includes sensing, computation and actuator delay.
Wet roads, worn tyres and downhill gradients reduce the available deceleration.

A safe system needs an operational design domain: the conditions in which automation is allowed. These
conditions may specify road type, speed, weather and map quality. At the boundary of this domain, the
vehicle must request human control or perform a minimal-risk manoeuvre. Safety therefore depends on
clear limits and predictable fallback behaviour, not only on intelligent algorithms.
\end{ecarticle}

\begin{tcolorbox}[title=\sffamily\bfseries Did You Know?,colback=yellow!10,colframe=ECOrange]
A vehicle may correctly detect an object yet still choose an unsafe action. Verification must therefore
cover the complete chain from sensing to decision and control.
\end{tcolorbox}

\begin{engineeringfocus}
\textbf{Autonomous driving pipeline}
\begin{center}
\resizebox{0.98\linewidth}{!}{%
\begin{tikzpicture}[font=\scriptsize,>=Latex,node distance=3mm]
\node[draw=ECBlue,fill=ECLightBlue,rounded corners] (sense) {sensors};
\node[draw=ECGreen,fill=ECGreen!10,rounded corners,right=of sense] (per) {perception};
\node[draw=ECPurple,fill=ECPurple!10,rounded corners,right=of per] (pred) {prediction};
\node[draw=ECOrange,fill=ECOrange!10,rounded corners,right=of pred] (plan) {planning};
\node[draw=ECRed,fill=ECRed!8,rounded corners,right=of plan] (ctrl) {control};
\node[draw=black,fill=gray!10,rounded corners,below=of plan] (safe) {safety monitor};
\draw[->,thick] (sense)--(per); \draw[->,thick] (per)--(pred); \draw[->,thick] (pred)--(plan); \draw[->,thick] (plan)--(ctrl);
\draw[->,thick,dashed] (safe)--(ctrl); \draw[->,thick,dashed] (per)--(safe);
\end{tikzpicture}}
\end{center}
\end{engineeringfocus}

\begin{keyfigures}
\begin{tabularx}{\linewidth}{@{}Xr@{}}
Camera strength & detailed images\\
Radar strength & relative speed\\
Lidar output & 3D point cloud\\
Critical principle & redundancy\\
Safe fallback & minimal-risk manoeuvre
\end{tabularx}
\end{keyfigures}

\begin{tcolorbox}[title=\sffamily\bfseries Illustrative braking distance,colback=white,colframe=ECOrange]
\begin{center}
\begin{tikzpicture}
\begin{axis}[width=0.96\linewidth,height=46mm,xlabel={Speed (km/h)},ylabel={Braking distance (m)},grid=major,
legend style={font=\scriptsize,at={(0.5,-0.22)},anchor=north,legend columns=2},ticklabel style={font=\scriptsize}]
\addplot[thick,mark=*] table[x=speed,y=dry,col sep=comma]{data/EC23/stopping.csv};
\addplot[thick,mark=square*] table[x=speed,y=wet,col sep=comma]{data/EC23/stopping.csv};
\legend{Dry road,Wet road}
\end{axis}
\end{tikzpicture}
\end{center}
\end{tcolorbox}

\begin{tcolorbox}[title=\sffamily\bfseries Quantitative application,colback=white,colframe=ECBlue]
A shuttle travels at $36\,\mathrm{km\,h^{-1}}$. The sensing and computation delay is $0.35\,\mathrm{s}$,
and the brake actuator delay is $0.15\,\mathrm{s}$. Assume a constant deceleration of
$5.0\,\mathrm{m\,s^{-2}}$.
\begin{enumerate}
\item Convert the speed into $\mathrm{m\,s^{-1}}$.
\item Calculate the distance travelled during the total delay.
\item Using $d_b=v^2/(2a)$, calculate the braking distance.
\item Determine the total stopping distance and propose a minimum detection range with a 25\% safety margin.
\end{enumerate}
\end{tcolorbox}

\begin{vocabulary}
\begin{tabularx}{\linewidth}{@{}lX@{}}
sensor fusion & fusion de capteurs\\
point cloud & nuage de points\\
localisation & localisation\\
trajectory & trajectoire\\
fallback & solution de repli\\
operational design domain & domaine opérationnel\\
wheel slip & glissement des roues\\
blind spot & angle mort\\
braking distance & distance de freinage\\
redundancy & redondance
\end{tabularx}
\end{vocabulary}

\begin{questions}
\item What functions form the autonomous driving chain?
\item Why are several sensor technologies combined?
\item How can localisation fail in a tunnel?
\item Why does braking distance rise rapidly with speed?
\item What is an operational design domain?
\end{questions}

\begin{engineeringchallenge}
Design an autonomous shuttle for a closed university campus. Define its sensors, maximum speed,
operational design domain, emergency strategy and verification tests for pedestrians, rain and GNSS
loss.
\end{engineeringchallenge}

\begin{speaking}
Who should be responsible after an autonomous vehicle accident: the passenger, manufacturer,
software provider or road operator? Defend a position.
\end{speaking}
\end{multicols}

\subsection*{Sources}
\ECsource{SAE International -- Automated Driving}{https://www.sae.org/standards/content/j3016_202104/}
\ECsource{NHTSA -- Automated Vehicles}{https://www.nhtsa.gov/vehicle-safety/automated-vehicles-safety}